\section{\demonstrate--\search--\predict} \label{section:dsp} \label{sec:dsp}

We now introduce the \DSP framework and show its expressive power by suggesting a number of strategies in which the \LM{} and \RM{} can come together to tackle complex problems effectively. We show in \secref{sec:eval} that such strategies outperform existing in-context learning methods. We begin by discussing the \LM{} and \RM{} foundation modules on which \DSP is built (\secref{sec:dsp:modules}) and then the datatypes and control flow within \DSP (\secref{sec:dsp:concepts}). Subsequently, we discuss each of the three inference stages: \demonstrate (\secref{sec:dsp:demonstrate}), \search (\secref{sec:dsp:search}), and \predict (\secref{sec:dsp:predict}).








\subsection{Pretrained Modules: \LM{} and \RM{}}
\label{sec:dsp:modules}

A \DSP program defines the communication between the language model \LM{} and the retrieval model \RM{}.%

\vspace{-2mm}
\paragraph{Language Model} We invoke a frozen language model \LM{} to conditionally \texttt{generate} (or \texttt{score}) text. For each invocation, the program prepares a \textit{prompt} that adapts the \LM{} to a specific function (e.g., answering questions or generating queries). A prompt often includes instructions, a few demonstrations of the desired behavior, and an input query to be answered.

As in Figure~\ref{fig:example}, the \LM{} generates not only: \textbf{(i)} the final answer to the input question (in the \predict stage), but also \textbf{(ii)} intermediate ``hop'' queries to find useful information for the input question (\search)  as well as \textbf{(iii)} exemplar queries that illustrate how to produce queries for questions in the training set (\demonstrate). This systematic use of the \LM{} is a hallmark of \DSP programs.


\vspace{-2mm}
\paragraph{Retrieval Model} \DSP programs also invoke a frozen retrieval model \RM{} to \texttt{retrieve} the top-$k$ most ``relevant'' text sequences for a given \textit{query}. The \RM{} can \texttt{index} a massive set of pre-defined \textit{passages} for scalable search, and those passages can be updated without changing the retrieval parameters. %
The \RM{} accepts free-form textual inputs and specializes in estimating the relevance (or similarity) of a text sequence to a query.

As in Figure~\ref{fig:example}, the \RM{} is responsible for retrieving \textbf{(i)} passages for each query generated by the \LM{} (during the \search stage), but also \textbf{(ii)} passages that are used within demonstrations (\demonstrate). In the latter case, the \RM's contributions are less about providing directly relevant information to the input question and more about helping the \LM{} adapt to the domain and task.

Though not utilized in this example, the \RM{} is also used in \DSP for functions like retrieving ``nearest-neighbor'' demonstrations from task training data (\demonstrate) and selecting well-grounded generated sequences from the \LM{} (\predict).

\subsection{Datatypes and Control Flow}
\label{sec:dsp:concepts}

We have implemented the \DSP{} framework in Python. The present section introduces the core data types and composable functions provided by the framework. We use illustrative code snippets to ground the examples, and to convey the power that comes from being able to express complex interactions between the \LM{} and \RM{} in simple programs.

\vspace{-2mm}
\paragraph{The \texttt{Example} Datatype} To conduct a task, a \DSP program manipulates one or more instances of the \texttt{Example} datatype. An \texttt{Example} behaves like a Python dictionary with multiple fields. The program is typically provided with a few training examples. The code snippet below illustrates this for multi-hop question answering.

\begin{samepage}
\begin{lstlisting}[language=Python,breaklines=true,showstringspaces=false]
from dsp import Example

train = [Example(question="When was the discoverer of Palomar 4 born?", answer="1889"),
         Example(question="In which city did Akeem Ellis play in 2017?", answer="Ellesmere Port")]
\end{lstlisting}
\end{samepage}

This snippet contains two labeled examples, each with a multi-hop question (e.g., ``In which city did Akeem Ellis play in 2017?'')\ and its short answer (``Ellesmere Port''). Arbitrary keys and values are allowed within an \texttt{Example}, though typical values are strings or lists of strings. 

In this task, we are unlikely to find an individual passage that provides the answer to any question. For example, the first training example can probably be resolved only by first answering the question of who discovered Palomar (``Edwin Hubble'') and then addressing the question of Hubble's birth date using different evidence passages. We typically assume that the human-labeled training data do \textit{not} include labels for intermediate transformations (e.g., queries for individual hops) that would be useful for following these steps, and so it is the job of the \DSP{} program to discover these strategies via in-context learning. %

\vspace{-2mm}
\paragraph{A \DSP Program} The following code snippet is a complete program for resolving multi-hop questions like those in Figure~\ref{fig:tasks}, with help from train examples like those above.

\begin{samepage}
\begin{lstlisting}[language=Python,breaklines=true,showstringspaces=false]
def multihop_program(question: str) -> str:
   x = Example(question=question, train=train)
   x = multihop_demonstrate(x)
   x = multihop_search(x)
   x = multihop_predict(x)
   return x.answer

multihop_program("How many storeys does the castle David Gregory inherited have?")
# => "five storeys"
\end{lstlisting}
\end{samepage}

The program takes the input (here, a question) and outputs the system output (its short answer). It starts by creating an \texttt{Example} for the input question and assigning the \texttt{train} field to the training set from the previous snippet. Programs invoke and compose \DSP \textit{primitives} (i.e., built-in functions) to build the \demonstrate, \search, and \predict transformations that define the program.

\vspace{-2mm}
\paragraph{Transformations} A transformation is a function that takes an \texttt{Example} as input and returns an  \texttt{Example}, populating new fields (or modifying existing fields) in it. 
This program invokes three developer-defined transformations, namely, \texttt{multihop\_demonstrate}, \texttt{multihop\_search}, and \texttt{multihop\_predict}. Transformations may themselves invoke other transformations, and they act analogously to layers in standard deep neural network (DNN) programming frameworks such as PyTorch, except that they pass text data instead of tensors between each other and do not involve backpropagation.

We categorize transformations according to their behavior (or purpose) under one of the \demonstrate, \search, and \predict stages. That said, \DSP does not impose this categorization and allows us to define functions that may blend these stages. We will discuss each of the three stages next.

\subsection{\demonstrate}
\label{sec:dsp:demonstrate}




 It is known that including examples of the desired behavior from the \LM{} in its prompt typically leads to better performance \cite{brown2020language}. In \DSP, a \textit{demonstration} is a training example that has been prepared to illustrate specific desired behaviors from the \LM{}. A \demonstrate transformation takes as input \texttt{x} of type \texttt{Example} and prepares a list of demonstrations in \texttt{x.demos}, typically by \textit{selecting} a subset of the training examples in \texttt{x.train} and \textit{bootstrapping} new fields in them.

\vspace{-2mm}
\paragraph{Bootstrapping Demonstrations} Examples in the training set typically consist of the input text and the target output of the task. The \demonstrate stage can augment a training example by programmatically bootstrapping annotations for intermediate transformations. In our running ``multi-hop'' example, the demonstrations illustrate three \LM{}-based transformations: \textbf{(i)} how to break down the input question in order to gather information for answering it (i.e., first-hop retrieval), \textbf{(ii)} how to use information gathered in an earlier ``hop'' to ask follow-up questions, and \textbf{(iii)} how to use the information gathered to answer complex questions.

\begin{samepage}
\begin{lstlisting}[language=Python,breaklines=true,showstringspaces=false]
Examples = list[Example]
Transformation = Callable[[Example],
                          Optional[Example]]

annotate(train: Examples, fn: Transformation)
    -> Examples
\end{lstlisting}
\end{samepage}


Akin to a specialized \texttt{map}, the \texttt{annotate} primitive accepts a user-defined transformation \texttt{fn} and applies it over a list of training examples. Whenever \texttt{fn} returns an example (rather than \texttt{None}), \texttt{annotate} caches the intermediate predictions (i.e., the generated queries and retrieved passages). These predictions serve as successful demonstrations for the pipeline transformations. In simple uses, \texttt{fn} may attempt to answer the example ``zero-shot'' one or more times. This is typically done by invoking the \search and \predict stages of the program. When an answer is produced, if \texttt{fn} assesses it as correct, it returns a populated example in which the intermediate predictions are present.

\vspace{-2mm}
\paragraph{Case Study}

The snippet below defines the function \texttt{multihop\_demonstrate}, called in Line~3 of \texttt{multihop\_program}, and illustrates the usage of \texttt{annotate}.

\begin{samepage}
\begin{lstlisting}[language=Python,breaklines=true]
from dsp import sample, annotate

def attempt_example(d: Example):
   d = d.copy(demos=[])
   d = multihop_search(d)
   d = multihop_predict(d)
   return d if d.pred == d.answer else None

def multihop_demonstrate(x: Example):
   demos = annotate(x.train, attempt_example)
   return Example(x, demos=demos)
\end{lstlisting}
\end{samepage}

In Line~10, \texttt{multihop\_demonstrate} invokes \texttt{annotate}, which bootstraps missing fields in training examples by caching annotations from \texttt{attempt\_example}. The transformation \texttt{attempt\_example} takes a training example \texttt{d} and attempts to answer it in a zero-shot fashion: it creates a copy of \texttt{d} with no demonstrations (Line~4; i.e., zero-shot) and invokes the multi-hop search and predict pipeline (Lines~5 and~6). Each transformation returns an updated version of \texttt{d} with additional fields populated. If the pipeline answers correctly (Line~7), the updated \texttt{d} is returned.

Figure~\ref{fig:example} illustrates this behavior. \demonstrate transforms a training question--answer pair to a fully-populated demonstration, including fields such as \texttt{hop1} and \texttt{hop2} (i.e., queries for multi-hop search) as well as \texttt{psg1} and \texttt{psg2}. When the \LM{} is later invoked to conduct a transformation, say, generating a ``second-hop'' query during \search, the \texttt{psg1} field serves as context and the \texttt{hop2} field serves as a label for this particular training example.


\vspace{-2mm}
\paragraph{Discussion}

This simple case study illustrates the power of composition in the \DSP abstraction. Because the pipeline is a well-defined program in which transformations communicate by passing text attached to \texttt{Example}s, a simple map-and-filter strategy can leverage the \LM{} and \RM{} to bootstrap annotations for a full \textit{pipeline} from end-task labels. This is an extensible strategy, but even in its simplest form it generalizes the approaches explored recently by \citet{zelikman2022star}, \citet{wei2022chain}, \citet{zhang2022automatic}, and \citet{huang2022large} in which an \LM{} self-generates chain-of-thought rationales for an individual prompt.


By bootstrapping pipelines, \demonstrate makes it easy to explore complex strategies in \search and \predict without writing examples for every transformation. This includes strategies that are challenging to explore without custom annotations in traditional retrieval-augmented NLP. For instance, \citet{khattab2021baleen} introduces a pipeline for multi-hop reasoning that is trained with weak supervision, extending work by \citet{lee2019latent} and \citet{khattab2021relevance}. In it, the target 3 or 4 passages that need to retrieved must be labeled but the system discovers the best \textit{order} of ``hops'' automatically.

In contrast, \DSP allows us to build complex pipelines without labels for intermediate steps, because we can compose programs out of small transformations. If \LM{} and \RM{} can accurately process such transformations ``zero-shot'' (i.e., without demonstrations) on at least one or two examples, these examples can be discovered with end-task labels and used as demonstrations.


To draw on our earlier analogy with DNN frameworks like PyTorch, \demonstrate aims to replace the function of backpropagation in extensible ways by simulating the behavior of the program (corresponding to a ``forward'' pass) and programmatically learning from errors. In doing this with frozen models and with only end-task labels, \demonstrate introduces a high degree of modularity. In particular, without hand-labeling intermediate transformations, developers may swap the training domain, update the training examples, or modify the program's strategy, and use \texttt{annotate} to automatically populate all of the intermediate fields for demonstrations.















\vspace{-2mm}
\paragraph{Selecting Demonstrations} It is not always possible to fit all of the training examples in the context window of the \LM{}. \DSP provides three primitives for selecting a subset of training examples, namely, \texttt{sample}, \texttt{knn}, and \texttt{crossval}.

\begin{samepage}


\begin{lstlisting}[language=Python,breaklines=true,showstringspaces=false]
sample(train: Examples, k: int)
    -> Examples

knn(train: Examples, cast: Callable[[Example], str])
    -> fn(example: Example, k: int) # currying
    -> Examples

crossval(train: Examples, n: int, k: int)
    -> fn(evaluate: Transformation)
    -> Examples
\end{lstlisting}
\end{samepage}

As a baseline choice, $k$ demonstrations can be randomly sampled from \texttt{train} using the \texttt{sample} primitive, an approach used by \citet{brown2020language} and much subsequent work. We can also leverage the \RM{}'s representations and select from the training set the $k$ nearest neighbors to the input text, a strategy explored by \citet{liu2021makes}. Another strategy is to apply cross-validation to select among a number of sampled sets of demonstrations \cite{perez2021true}. For example, given $|\texttt{train}| = 100$ training examples, \texttt{crossval} would select $n$ subsets of $k = 5$ examples each, and return the set with which a transformation \texttt{evaluate} performs best on the remaining $95$ examples.

\vspace{-2mm}
\paragraph{Compositions \& Extensions}

By manipulating demonstrations and higher-order transformations, these simple selection and bootstrapping primitives can be combined to conduct larger novel strategies. If the training set is very large (e.g., $|\texttt{train}| = 100,000$), we can conduct \texttt{knn} to find the nearest $k=16$ examples and only \texttt{annotate} these, arriving at a system that learns incrementally in real-time. If the training set is moderately large (e.g., $|\texttt{train}| = 1000$), we can conduct \texttt{crossval} and cache the performance of all prompts it evaluates on each training example. At test time, we can use \texttt{knn} to find $k=50$ similar examples to the test input and select the prompt that performs best on these $k$ examples, producing an adaptive system that is informed by the quality of its pipeline on different types of examples.


 


\subsection{\search}
\label{sec:dsp:search}








The \search stage gathers passages to support transformations conducted by the \LM{}. We assume a large knowledge corpus---e.g., a snippet of Web, Wikipedia, or arXiv---that is divided into text \textit{passages}.
Providing passages to the \LM{} facilitates factual responses, enables updating the knowledge store without retraining, and presents a transparency contract: when in doubt, users can check whether the system has faithfully used a reliable source in making a prediction.


In the simplest scenarios, \search can directly query the \RM{}, requesting the top-$k$ passages (from a pre-defined index) that match an input question. This baseline instantiation of \search simulates retrieval in most open-domain question answering systems, which implement a ``retrieve-then-read'' pipeline, like \citet{lee2019latent}, \citet{khattab2021relevance}, \citet{lazaridou2022internet}, and many others.

\begin{samepage}

\begin{lstlisting}[language=Python,breaklines=true]
from dsp import retrieve

def simple_search(x):
    passages = retrieve(query=x.question, k=2)
    return passages
\end{lstlisting}
\end{samepage}

\vspace{-2mm}
\paragraph{\search Strategies}



In many scenarios, the complexity of the task demands more sophisticated \search strategies that empower the \RM{} to find relevant passages. Our running example (Figure~\ref{fig:example}) is one such scenario, in which we suspect examples are likely to require \textit{multi-hop} reasoning in particular. Other settings, for instance, pose conversational challenges, in which the information need expressed by a user can only be resolved by taking into account previous turns in the conversation, or demand more extensive planning~\cite{zhong2022romqa}.

In the retrieval-augmented NLP literature, multi-hop search~\cite{xiong2020answering,khattab2021baleen} and conversational search~\cite{del2021question,raposo2022question} pipelines have received much attention. These systems are typically fine-tuned with many hand-labeled query ``rewrites''~\cite{anantha2020open}, ``decompositions''~\cite{geva2021did,min2019multi}, or target hops~\cite{yang2018hotpotqa,jiang2020hover}. %
Supported with automatic annotations from \demonstrate, the \search stage allows us to simulate many such strategies and many others
in terms of passing queries, passages, and demonstrations between the \RM{} and \LM{}. More importantly, \search facilitates our vision of advanced strategies in which the \LM{} and \RM{} cooperate to incrementally plan a research path for which the \RM{} gathers information and the \LM{} identifies next steps.

\paragraph{Case Study} Let us build on our running multi-hop example as a case study. We can define \texttt{multihop\_search\_v2} (Line~4 in our core program), a slightly more advanced version of the \search transformation from Figure~\ref{fig:example}. This transformation simulates the iterative retrieval component of fine-tuned retrieval-augmented systems like IRRR~\cite{qi2020retrieve}, which reads a retrieved passage in every hop and generates a search query (or a termination condition to stop hopping), and Baleen~\cite{khattab2021baleen}, which summarizes the information from many passages in each hop for inclusion in subsequent hops.

\begin{samepage}

\begin{lstlisting}[language=Python]
from dsp import generate

def multihop_search_v2(x, max_hops=3):
  x.hops = []
  
  for hop in range(max_hops):
    summary, query = generate(hop_template)(x)
    x.hops.append((summary, query))
    
    if query == 'N/A': break
    
    passages = retrieve(query, k=5)
    x.context = [summary] + passages
    
  return x
\end{lstlisting}
\end{samepage}




In \texttt{multihop\_search\_v2}, Line~7 calls the \texttt{generate} primitive, which invokes the \LM{} to produce a query for each retrieval hop. The \LM{} is conditioned on a prompt that is prepared using the \texttt{hop\_template} template. (We discuss prompt templates and the \texttt{generate} primitive in~\secref{sec:dsp:predict}.) Here, this template may be designed to generate a prompt that has the following format (e.g., for the second hop).

\begin{samepage}
\begin{lstlisting}[breaklines=true]
My task is to write a simple query that gathers information for answering a complex question. I write N/A if the context contains all information required.

{Task demonstrations from x.demos, if any}

Context: {x.context}
Question: {x.question}
Summary: Let's summarize the above context. __{summary}__
Search Query: __{query}__
\end{lstlisting}
\end{samepage}

As shown, the \LM{} is instructed to read the context retrieved in earlier hops and a complex question. It is then prompted to write: \textbf{(i)} a summary of the supplied context and \textbf{(ii)} a search query that gathers information for answering that question. The generated text will be extracted and assigned to the \texttt{summary} and \texttt{query} variables in (\texttt{multihop\_search\_v2}; Line~7). On Line~10, we terminate the hops if the query is ``N/A''. Otherwise, Line~12 retrieves $k=5$ passages using the query and Line~13 assigns the \texttt{context} for the subsequent hop (or for \predict), setting that to include the \texttt{summary} of all previous hops as well as the passages retrieved in the final hop so far.

\paragraph{Comparison with self-ask} It may be instructive to contrast this multi-hop \DSP program with the recent ``self-ask''~\cite{press2022measuring} prompting technique, which we compare against in~\secref{sec:eval}. Self-ask can be thought of as a simple instantiation of \DSP's \search stage. In it, the \LM{} asks one or more ``follow-up questions'', which are intercepted and sent to a search engine. The search engine's answers are concatenated into the prompt and are used to answer the question. This is essentially a simplified simulation of IRRR~\cite{qi2020retrieve}.

As a general framework, \DSP can express ideas like self-ask and many other, more sophisticated pipelines as we discuss in the present section. More importantly, \DSP offers a number of intrinsic advantages that lead to large empirical gains: 80\%--290\% over self-ask. For instance, \DSP programs are deeply modular, which among other things means that \DSP programs will annotate and construct their own demonstrations. Thus, they can be developed without labeling any of the intermediate transformations (e.g., the queries generated). In addition, the \LM{} prompts constructed by \DSP get automatically updated to align with the training data and retrieval corpus provided. In contrast, approaches like self-ask rely on a hand-written prompt with hard-coded examples.

Moreover, \DSP assigns the control flow to an explicit program and facilitates design patterns that invoke the \LM{} (or \RM{}) to conduct small transformations. This allows us to build steps that are dedicated to generating one or more retrieval queries, summarizing multiple passages per hop, and answering questions. These steps are individually simpler than the self-ask prompt, yet our multi-hop \DSP program deliberately composes them to build richer pipelines that are thus more reliable. In contrast, self-ask delegates the control flow to the \LM{} completions, maintaining state within the prompt itself and intercepting follow-up questions to conduct search. We find that this paradigm leads to a ``self-distraction'' problem (\secref{sec:eval:datasets}) that \DSP programs avoid.


\paragraph{Fusing Retrieval Results} For improved recall and robustness, we can also \textit{fuse} the retrieval across multiple generated queries. Fusion has a long history in information retrieval~\cite{fox1994combination,xue2013modeling,kurland2018fusion} and sequentially processing multiple queries was explored recently by \citet{gao2022attributed} for retroactively attributing text generated by LMs to citations. Inspired by these, we include a \texttt{fused\_retrieval} primitive to \DSP to offer a versatile mechanism for interacting with frozen retrievers. It accepts an optional fusion function that maps multiple retrieval lists into one. By default, \DSP uses a variant of CombSUM~\cite{fox1994combination}, assigning each passage the sum of its probabilities across retrieval lists.

To illustrate, the modification below generates $n=10$ queries for the transformation \texttt{multihop\_search\_v2}.

\begin{samepage}
\begin{lstlisting}[numbers=none,language=Python,breaklines=true]
c = generate(hop_template, n=10)(x)
passages = fused_retrieval(c.queries, k=5)
summary = c.summaries[0] # highest-scoring summary
\end{lstlisting}
\end{samepage}

\paragraph{Compositions \& Extensions} To illustrate a simple composition, we can equip a chatbot with the capacity for conversational multi-hop search by combining a query rewriting step, which produces a query that encompasses all of the relevant conversational context, with the multi-hop transformation, as follows.

\begin{samepage}
\begin{lstlisting}[language=Python,breaklines=true]
def conversational_multihop_search(x):
    x.question = generate(conv_rewriting_template)(x)
    return multihop_search_v2(x)
\end{lstlisting}
\end{samepage}

Similar approaches can be used for correcting spelling mistakes or implementing pseudo-relevance feedback~\cite{cao2008selecting,wang2022colbert}, in which retrieved passages are used to inform a better search query, though this has not been attempted with pretrained LMs to our knowledge.







\subsection{\predict}
\label{sec:dsp:predict}

The \predict stage generates the system output using demonstrations (e.g., in \texttt{x.demos}) and passages (e.g., in \texttt{x.context}). \predict tackles the challenges of reliably solving the downstream task, which integrates much of the work on in-context learning in general. Within \DSP, it also has the more specialized function of systematically aggregating information across a large number of demonstrations, passages, and candidate predictions.

\paragraph{Generating Candidates} Generally, \predict has to produce one or more candidate predictions for the end-task. To this end, the basic primitive in \predict is \texttt{generate}, which accepts a \texttt{Template} and (via currying) an \texttt{Example} and queries the \LM{} to produce one or more completions, as explored earlier in~\secref{sec:dsp:search}. A corresponding primitive that uses the \RM{} in this stage is \texttt{rank}, which accepts a query and one or more passages and returns their relevance scores.

\begin{samepage}

\begin{lstlisting}[language=Python,breaklines=true]
Template  # template: an object that can produce prompts and parse completions

generate(template: Template)
    -> fn(example: Example)
    -> Completions  # object with keys to access extracted preds and scores

rank(query: str, passages: List[str])
    -> List[float]  # object with keys to access passage texts and scores
\end{lstlisting}
\end{samepage}



A \texttt{Template} is an object that can produce prompts, that is, map an \texttt{Example} to a string, and extract fields out of completions. For instance, we can map an example \texttt{x} that has a question and retrieved passages to the following prompt:

\begin{samepage}

\begin{lstlisting}[breaklines=true]
My task is to answer questions using Web documents.

{Task demonstrations from x.demos, if any}

Context: {x.passage}
Question: {x.question}
Rationale: Let's think step by step. __{rationale}__
Answer: __{answer}__
\end{lstlisting}
\end{samepage}

As this illustrates, the \LM{} will be asked to generate a chain-of-thought rationale (CoT; \citealt{wei2022chain}; \citealt{kojima2022large}) and an answer, and the generated text will be extracted back into the \texttt{rationale} and \texttt{answer} keys of each completion.

Each invocation to the \LM{} can sample multiple candidate predictions. Selecting a ``best'' prediction is the subject of much work on decoding \cite{wiher2022decoding,li2022contrastive}, but a frozen and general-purpose \LM{} may not support custom modifications to decoding. Within these constraints, we present several high-level strategies for selecting predictions and aggregating information in \DSP via the \LM{} and \RM{}. 









\paragraph{Selecting Predictions} Among multiple candidates, we can simply extract the most popular prediction. When a CoT is used to arrive at the answer, this is the self-consistency method of \citet{wang2022self}, which seeks to identify predictions at which multiple distinct rationales arrive.

\begin{samepage}
\begin{lstlisting}[language=Python,breaklines=true]
from dsp import generate, majority

def multihop_predict(x):
   candidates = generate(template_qa)(x)
   return x.copy(answer=majority(candidates).answer)
\end{lstlisting}
\end{samepage}

\DSP generalizes this in two ways. First, we can sample multiple ``pipelines of transformations'' (PoT) within the program, rather than locally with ``chains of thought'' (CoT) in one transformation. These chains may even invoke different paths in the program, as illustrated below.

\vspace{8mm}
\begin{samepage}
\begin{lstlisting}[language=Python,breaklines=true]
from dsp import branch

def pipeline(x):
  return multihop_predict(multihop_search_v2(x))

def PoT_program(question: str) -> str:
  x = Example(question=question, train=train)
  x = multihop_demonstrate(x)
  
  candidates = branch(pipeline, n=5, t=0.7)(x)
  return x.copy(answer=majority(candidates).answer)
\end{lstlisting}
\end{samepage}


In the snippet above, Line~10 invokes the primitive \texttt{branch} which samples $n$ different PoTs with a high temperature (e.g., $t=0.7$) and accumulates their intermediate and final predictions. In this example, our pipeline invokes \texttt{multihop\_search\_v2} (\secref{sec:dsp:search}), which applies a variable number of retrieval hops depending on the questions generated, before doing \predict. That is, \texttt{PoT\_program} potentially invokes multiple distinct paths in the program (i.e., with different multi-hop queries and number of hops in each) across branches. It then selects the \texttt{majority} answer overall.

\DSP generalizes self-consistency in a second way. When sampling our CoTs or PoTs provides multiple candidates, we can select the top-$k$ (e.g., top-4) predictions and then \textit{compare} them directly. For instance, we may prompt the \LM{} to compare these choices as MCQ candidates, a transformation for which \demonstrate can automatically prepare exemplars. This effectively simulates the LM recursion of \citet{levine2022standing}, though unlike their approach it does not require a large training set or updating any (prompt-tuning) weights. One such implementation is illustrated in \texttt{openqa\_predict} below.

\begin{samepage}
\begin{lstlisting}[language=Python,breaklines=true]
def openqa_predict(x):
   preds = generate(template_qa, n=20)(x).answers
   x.choices = most_common(preds, k=4)
   
   queries = [f"{x.question} {c}"
              for c in x.choices]
              
   x.passages = fused_retrieval(queries)
   x.answer = generate(TemplateMCQ)(x).answer
   return x
\end{lstlisting}
\end{samepage}

As an alternative comparison approach, we can invoke the \RM{} via \texttt{rank} to find the prediction that is most grounded in a retrieved contexts (i.e., most similar to the concatenation of the retrieved passages) or, given an \RM{} that can score completions~\cite{krishna2022rankgen}, simply the prediction that has the highest score given the prompt.

\vspace{-2mm}
\paragraph{Aggregating Information}  When only a few demonstrations or passages are selected, we can simply concatenate them all into the prompt. For instance, GPT-3.5 \texttt{text-davinci-002} has a context window of 4097 tokens, which we find to be reasonably large for accommodating several (e.g., 3--5) demonstrations, which individually include their own passages and rationales.

To deal with a larger number of demonstrations or passages, we can \texttt{branch} in parallel to process individual subsets of the passages or demonstrations and then aggregate the individual answers using one of the scoring methods presented earlier. Indeed, \citet{lewis2020retrieval} and \citet{lazaridou2022internet} have explored marginalization as a way to combine scores across passages and \citet{le2022few} ensemble prompts across demonstrations, which can be expressed in this way.


An alternative aggregation strategy is to accumulate information across passages sequentially, rather than independently. This is effectively how our multi-hop approach works (\secref{sec:dsp:search}). Such a strategy has also been employed recently by \citet{gao2022attributed} for retroactively attributing text generated by LMs to citations. They generate many queries but instead of fusion (\secref{sec:dsp:search}), they run their pipeline on each query and use its outputs to alter the input to subsequent queries.\footnote{Though most of the functionality in this section is implemented, the primitives \texttt{branch}, \texttt{knn}, and \texttt{crossval} are currently work-in-progress.}






 
























